\chapter{一致可积性}

\begin{solution}{ui-equivalent-truncation}
在 \(\{|f|>2M\}\) 上，
\((|f|-M)_+\ge|f|/2\)，故
\[
 \int_{\{|f|>2M\}}|f|
 \le2\int(|f|-M)_+.
\]
反过来，
\((|f|-M)_+\le|f|\mathbf1_{\{|f|>M\}}\)。两边分别取函数族上确界并令
\(M\to\infty\)，得到等价性。
\end{solution}

\begin{solution}{ui-finite-family}
对每个 \(f_j\in L^1\)，
\(\int_{\{|f_j|>M\}}|f_j|\to0\)。有限多个趋零量的最大值仍趋零，所以
\[
 \sup_{1\le j\le r}\int_{\{|f_j|>M\}}|f_j|\to0.
\]
\end{solution}

\begin{solution}{ui-dominated-family}
若 \(|f|\le g\)，则
\[
 \int_E|f|\le\int_Eg.
\]
由 \(g\) 的积分绝对连续性，右侧在 \(\mu(E)\) 足够小时统一小。函数族还
满足 \(\sup_f\|f\|_1\le\|g\|_1=:C\)。给定 \(\varepsilon>0\)，
先由积分绝对连续性取 \(\delta>0\)，使 \(\mu(E)<\delta\) 时
\(\int_Eg<\varepsilon/2\)。取 \(M>(C+1)/\delta\)，Markov 不等式给出
\(\mu(\{|f|>M\})\le C/M<\delta\)。将该集合代入刚才的估计，
便有 \(\sup_f\int_{\{|f|>M\}}|f|<\varepsilon\)。这里没有用控制函数
自身的尾积分，也没有要求总测度有限。
\end{solution}

\begin{solution}{ui-spike-details}
\[
 \|n\mathbf1_{(0,1/n)}\|_1=1.
\]
对任意 \(\varepsilon>0\)，当 \(n>\varepsilon\) 时超水平集长度为
\(1/n\to0\)，故依测度趋零。若 \(M>0\)，取 \(n>M\)，则
\[
 \int_{\{|f_n|>M\}}|f_n|=1,
\]
所以尾积分上确界不趋零。
\end{solution}

\begin{solution}{ui-lp-bounded}
在 \(\{|f|>M\}\) 上，
\(|f|\le M^{1-p}|f|^p\)。因此
\[
 \int_{\{|f|>M\}}|f|
 \le M^{1-p}\int|f|^p
 \le C^pM^{1-p}\to0.
\]
\end{solution}

\begin{solution}{ui-convex-sufficient}
记 \(C=\sup_f\int\Phi(|f|)\)。由
\(\Phi(t)/t\to\infty\)，对给定 \(\varepsilon>0\) 可取 \(M\) 使
\[
 t\mathbf1_{\{t>M\}}\le\frac{\varepsilon}{C+1}\Phi(t).
\]
积分并取上确界即得尾积分至多为 \(\varepsilon C/(C+1)<\varepsilon\)，
包括 \(C=0\) 的情形。这个方向只用
\(\Phi(t)/t\to\infty\) 和非负性，没有用凸性；凸性用于反向构造和后续
函数空间理论中的良好性质。
\end{solution}

\begin{solution}{vitali-ae-version}
有限测度空间上 a.e. 收敛蕴含依测度收敛。其余条件正是 Vitali 定理的
假设，所以 \(\|f_n-f\|_1\to0\)。也可直接用 Egorov：删去小集合后
一致收敛，在补集上用有限测度控制积分，在小集合上用一致可积性控制。
\end{solution}

\begin{solution}{l1-convergence-ui}
给定 \(\varepsilon>0\)，取 \(N\) 使 \(n\ge N\) 时
\(\|f_n-f\|_1<\varepsilon/3\)。有限族
\(f,f_1,\ldots,f_{N-1}\) 的积分绝对连续性给出 \(\delta\)，使
\(\mu(E)<\delta\) 时每个相应积分小于 \(\varepsilon/3\)。于是对
\(n\ge N\)，
\[
 \int_E|f_n|
 \le\|f_n-f\|_1+\int_E|f|<2\varepsilon/3.
\]
有限初项也满足相同的小集合估计，极限 \(f\) 亦然，所以整个族满足
\textup{(UI-set)}。另取固定指标 \(N_0\)，使 \(n\ge N_0\) 时
\(\|f_n-f\|_1<1\)，则
\[
 C=\max\left\{\|f\|_1+1,\ \|f_1\|_1,\ldots,\|f_{N_0-1}\|_1\right\}
\]
是整个族的有限 \(L^1\) 上界；当 \(N_0=1\) 时只保留第一项。
由 Markov 不等式，任意族中函数 \(h\) 都满足
\(\mu(\{|h|>M\})\le C/M\)。给定小集合估计所需的 \(\delta\)，
取 \(M>(C+1)/\delta\) 并代入 \textup{(UI-set)}，得到尾积分统一趋零。
此证明适用于任意测度空间。
\end{solution}

\begin{solution}{ui-sums}
若 \(|f+g|>2M\)，则 \(|f|>M\) 或 \(|g|>M\)。并且
\[
 |f+g|\mathbf1_{\{|f+g|>2M\}}
 \le2|f|\mathbf1_{\{|f|>M\}}
   +2|g|\mathbf1_{\{|g|>M\}}.
\]
积分并分别取 \(\mathcal F,\mathcal G\) 上确界，右侧随
\(M\to\infty\) 越来越小。
\end{solution}

\begin{solution}{mass-tightness-translation}
因所有函数值不超过 \(1\)，当 \(M>1\) 时尾积分为零，故一致可积。设
\(K\subset\mathbb R\) 有限测度。若存在 \(\varepsilon<1\) 使对所有
\(n\) 都有
\(\int_{\mathbb R\setminus K}\mathbf1_{[n,n+1]}<\varepsilon\)，则
\(m(K\cap[n,n+1])>1-\varepsilon\) 对所有 \(n\) 成立。对互不相交的
半开区间 \([n,n+1)\) 求和会推出 \(m(K)=\infty\)，矛盾；改为半开区间
只删去零测端点，不改变积分。因此不具质量紧性。
\end{solution}

\begin{solution}{vitali-infinite-space}
取有限测度集 \(X_j\uparrow X\)。由局部依测度收敛，可递归选取严格递增的
\(n_j\)，满足
\[
 \mu\bigl(X_j\cap\{|f_{n_j}-f|>2^{-j}\}\bigr)<2^{-j}.
\]
坏集测度之和有限，所以几乎每个点只落入有限多个坏集；又每个点最终属于
\(X_j\)，故 \(f_{n_j}\to f\) a.e. 先由质量紧性取一个有限测度集 \(K_0\)，
使全体 \(f_n\) 在 \(X\setminus K_0\) 上的积分小于 \(1\)。有限测度 Vitali
定理说明 \(f\in L^1(K_0)\)，Fatou 引理说明
\(\int_{X\setminus K_0}|f|\le1\)，故 \(f\in L^1(X)\)。

现给定 \(\varepsilon>0\)，再由质量紧性选有限测度 \(K\)，使
\[
 \sup_n\int_{X\setminus K}|f_n|<\varepsilon.
\]
对上述 a.e. 收敛子列使用 Fatou，可得
\(\int_{X\setminus K}|f|\le\varepsilon\)。在 \(K\) 上，限制后的函数族
一致可积且依测度收敛，有限测度 Vitali 定理给
\(\int_K|f_n-f|\to0\)。因此
\[
 \limsup_n\|f_n-f\|_1\le2\varepsilon.
\]
令 \(\varepsilon\downarrow0\) 即得结论。尖峰列缺少一致可积性；移动指示
函数列缺少质量紧性；若没有局部依测度收敛，则即使另两项成立也不能确定
候选极限。
\end{solution}

\begin{solution}{ui-random-variables}
在概率空间 \((\Omega,\mathcal A,\mathbb P)\) 上，函数族
\(\mathcal H\subset L^1(\mathbb P)\) 的一致可积条件为
\[
 \lim_{M\to\infty}\sup_{H\in\mathcal H}
 \mathbb E\bigl(|H|\mathbf1_{\{|H|>M\}}\bigr)=0.
\]
题设已包含每个 \(|X_n|^p\) 可积，故 \(X_n\in L^p\)。总测度为1，
有限测度判据又给出 \(C=\sup_n\mathbb E|X_n|^p<\infty\)。
由子列原理取 \(X_{n_k}\to X\) 几乎处处，Fatou 引理给出
\(\mathbb E|X|^p\le C\)，故 \(X\in L^p\)。

令
\[
 Y_n=|X_n-X|^p.
\]
对每个 \(\varepsilon>0\)，有
\(\mathbb P(Y_n>\varepsilon)=\mathbb P(|X_n-X|>\varepsilon^{1/p})\to0\)，
所以 \(Y_n\to0\) 依概率。利用
\[
 |X_n-X|^p\le 2^{p-1}(|X_n|^p+|X|^p)
\quad(p\ge1)
\]
以及 \(0<p<1\) 时的
\(|a-b|^p\le|a|^p+|b|^p\)，得
\(Y_n\le c_p(|X_n|^p+|X|^p)\)，其中 \(c_p=\max\{1,2^{p-1}\}\)。
后一幂不等式由 \((u+v)^p\le u^p+v^p\) 得到：固定 \(v\ge0\)，
函数 \((u+v)^p-u^p\) 在 \(u>0\) 上递减，故不超过其在零点的值 \(v^p\)。

单个可积函数构成一致可积族，前面的和族习题说明
\(\{|X_n|^p+|X|^p\}\) 一致可积。若非负族 \(\{H_n\}\) 一致可积且
\(0\le Z_n\le cH_n\)，其中 \(c>0\) 固定，则
\[
 \mathbb E\bigl(Z_n\mathbf1_{\{Z_n>M\}}\bigr)
 \le c\,\mathbb E\bigl(H_n\mathbf1_{\{H_n>M/c\}}\bigr).
\]
取上确界并令 \(M\to\infty\)，可见 \(\{Z_n\}\) 也一致可积。
应用于 \(Y_n\)，再用有限测度 Vitali 定理，便得
\(\mathbb E|X_n-X|^p\to0\)。
\end{solution}
